\begin{abstract}
Some researchers speculate that intelligent reinforcement learning ({\rl}) agents would be incentivized to seek resources and power in pursuit of the objectives we specify for them. Other researchers point out that {\rl} agents need not have human-like power-seeking instincts. To clarify this discussion, we develop the first formal theory of the statistical tendencies of optimal policies. In the context of Markov decision processes ({\mdp}s), we prove that certain environmental symmetries are sufficient for optimal policies to tend to seek power over the environment. These symmetries exist in many environments in which the agent can be shut down or destroyed. We prove that in these environments, most reward functions make it optimal to seek power by keeping a range of options available and, when maximizing average reward, by navigating towards larger sets of potential terminal states.
\end{abstract}

\section{Introduction} 
\citet{omohundro_basic_2008,bostrom_superintelligence_2014,russell_human_2019} hypothesize that highly intelligent agents tend to seek power in pursuit of their goals. Such power-seeking agents might gain power over humans. Marvin Minsky imagined that an agent tasked with proving the Riemann hypothesis might rationally turn the planet—along with everyone on it—into computational resources \citep{russell_artificial_2009}. However, another possibility is that such concerns simply arise from the anthropomorphization of AI systems \citep{lecun_dont_2019,instrumental,steven_2020,mitchell2021ai}.

We clarify this discussion by grounding the claim that highly intelligent agents will tend to seek power. In \cref{sec:action-prob}, we identify optimal policies as a reasonable formalization of ``highly intelligent agents.''\footnote{This paper assumes that reward functions reasonably describe a trained agent's goals. Sometimes this is roughly true (\eg{} chess with a sparse victory reward signal) and sometimes it is not true. \citet{rewardNotOpt} argues that capable {\rl} algorithms do not necessarily train policy networks which are best understood as optimizing the reward function itself. Rather, they point out that—especially in policy gradient approaches—reward provides gradients to the network and thereby modifies the network's generalization properties, but doesn't ensure the agent generalizes to ``robustly optimizing reward'' off of the training distribution.} Optimal policies ``tend to'' take an action when the action is optimal for most reward functions. We expect future work to translate our theory from optimal policies to learned, real-world policies. 

\Cref{sec:power} defines ``power'' as the ability to achieve a wide range of goals. For example, ``money is power,'' and money is instrumentally useful for many goals. Conversely, it's harder to pursue most goals when physically restrained, and so a physically restrained person has little power. An action ``seeks power'' if it leads to states where the agent has higher power. 

We make no claims about when large-scale AI power-seeking behavior could become plausible. Instead, we consider the theoretical consequences of optimal action in {\mdp}s. \Cref{sec:symmetries} shows that power-seeking tendencies arise not from anthropomorphism, but from certain graphical symmetries present in many {\mdp}s. These symmetries automatically occur in many environments where the agent can be shut down or destroyed, yielding broad applicability of our main result (\cref{rsdIC}).